\cleardoublepage
\phantomsection
\addcontentsline{toc}{chapter}{DAFTAR SIMBOL}
\chapter*{DAFTAR SIMBOL}
\label{chap:simbol}
\setcounter{chapter}{0}
\setcounter{section}{0}

% Tugas Akhir ini menghasilkan artefak rekayasa perangkat lunak yang dievaluasi
% melalui metrik operasional, bukan model matematis dengan notasi khusus. Kuantitas
% yang dilaporkan (PSI, KS, p99 latensi, throughput, error budget) merupakan singkatan
% terminologis yang sudah dijelaskan pada Daftar Singkatan dan dijabarkan dalam teks
% pada bab perancangan, implementasi, dan evaluasi. Daftar di bawah hanya merangkum
% notasi non-singkatan yang muncul dalam teks.

\begin{longtable}{@{} l >{\raggedright\arraybackslash}p{8cm} c @{}}
\multicolumn{1}{l}{\textbf{Notasi}} & \multicolumn{1}{l}{\textbf{Deskripsi}} & \multicolumn{1}{l}{\parbox[c]{2.5cm}{\raggedright \textbf{Pemakaian pertama kali}}} \\
\hline

p99 & Persentil ke-99 dari distribusi latensi permintaan, dipakai sebagai \textit{indicator} \textit{tail latency} pada SLO layanan inferensi. & Bab~\ref{chap:perancangan} \\

\end{longtable}
